\documentclass[../main.tex]{subfiles}

\begin{document}

\chapter{DPLL}

The Davis-Putnam-Logemann-Loveland (DPLL) algorithm is a complete
backtracking based search algorithm for finding satisfiable assignments for SAT
problems. It is, dare I say, the foundational algorithm for modern SAT solvers,
as the state-of-the-art algorithm used by basically every SAT solver builds upon
the foundations that DPLL provides, as we will see later on.

\noindent It consists of three components:

\begin{itemize}
    \item Recursive backtracking
    \item Unit propagation
    \item Pure literal elimination
\end{itemize}
How does it work?

Imagine a binary tree of variables, each with two branches representing their
assignment, one for true and one for false. Suppose we don't know anything
about our CNF equation. That is, from looking at it, any variable could have
any value, and we don't know how each assignment helps us in finding out
whether or not the equation is satisfiable. How can we figure out if we can
satisfy the equation without just brute forcing it?

The answer is \sout{making an ass out of you and me} to make assumptions.

\begin{figure}[H]
    \centering
    \begin{tikzpicture}
        \tikzset{edge from parent/.style={draw,edge from parent 
            path={(\tikzparentnode.south)-- +(0,-8pt)-| (\tikzchildnode)}}}
        \Tree [.a
                    [.b
                        [.c
                            [.SAT ]
                            [.UNSAT ]
                        ]
                        [.c
                            [.SAT ]
                            [.UNSAT ]
                        ] 
                    ]
                    [.b
                        [.c
                            [.SAT ]
                            [.SAT ]
                        ]
                        [.c
                            [.UNSAT ]
                            [.UNSAT ] 
                        ] 
                    ] 
                ]
    \end{tikzpicture}
    \caption{Example DPLL decision tree (the outcomes may not correspond to an
    actual CNF instance, sue me).}
\end{figure}

Going back to our hypothetical, we don't know anything about anything. We
could have a giant decision tree to explore with an exponential number of leaves
to evaluate. How does making assumptions help us? Do we just run a random number generator 
and if it returns 1, we go ``okay, it's satisfiable,'' and false otherwise?

No. We make assumptions about our literal assignments.Unlike brute forcing where we try 
every possible combination of variable assignments at once, in DPLL we start with a literal, 
we'll call it $A$. Suppose there is nothing in the equation that indicates whether or not 
$A$ should be true or false; either could be valid.

So to tackle our problem, initially we will assume $A$ is true and see where
that leads us. We propagate that value throughout our clauses, and suppose
there is a clause that can't be satisfied. That's a contradiction, because if the
equation was satisfiable then $A = True$ couldn't have led us there.

So we've figured out that at the very least $A = True$ leads to a
contradiction. We've now eliminated half of the search space immediately and only need to
consider the cases where $A = False$.

This may not make total sense immediately, so instead of giving one explanation
of the algorithm and then building it, we will build up the algorithm and explain
how each part works.

\footnote{Note for future me: write a better introduction.}

\footnote{\url{https://en.wikipedia.org/wiki/DPLL_algorithm}}

\section{Recursive Backtracking}

First, we need to rewrite our brute force algorithm to be recursive, as that's
the foundation for the entire algorithm (duh). To do this, we will introduce a
new method \textit{solve()} which will be our recursive solving method. To
also make things simpler, we will extract the equation evaluation into its own
separate method as well.

This means our SAT solver becomes like so:

\begin{lstlisting}
public interface:
    SATSolver(string input)
    bool solveCNF()

private interface:
    void parse(string equation)
    bool solve(bool[] assignment)
    bool evaluate(bool[] assignment)
    int numVars
    int numClauses
    int[][] clauses
    optional bool[] lits
\end{lstlisting}

The flow essentially becomes like this: \textit{solveCNF()} calls
\textit{solve()}, which calls itself over and over and returns whatever solution
it finds back up to \textit{solveCNF()}, who will handle any sort of cleanup.

Why can't we just make \textit{solveCNF()} call itself recursively? Good question.
That's because \textit{solveCNF()} will handle any setup that
\textit{solve()} will need in order to function properly, and it will handle any
cleanup that \textit{solve()} may create.

Let's now discuss how \textit{solve()} will work.
It will be passed in a \textit{bool[]} for the current assignment of literals.
When the size of our assignment is equal to the number of variables, we will
then evaluate the satisfiability of the assignment, and if it's satisfiable we
return true. Otherwise, we return false, backtrack, retry, until we either find a
satisfiable assignment or we've exhausted them all.

The pseudocode is as so, starting with \textit{evaluate()}:

\begin{lstlisting}
bool evaluate(bool[] assignment)
{
    bool solved = true;

    for each (clause in clauses)
    {
        bool satisfied_clause = false;

        for (j = 1 up to clause.size)
        {
            int variable = clause[j];
            int v_idx = abs(variable);

            if (variable > 0)
            {
                satisfied_clause = assignment[v_idx]
                  || satisfied_clause;
            }
            else
            {
                satisfied_clause = !assignment[v_idx] 
                  || satisfied_clause;
            }

            if (satisfied_clause == True)
            {
                continue;
            }
        }

        solved &= satisfied_clause;

        if (solved == False)
        {
            break;
        }
    }

    if (solved == True)
    {
        lits = assignment;
    }
    else
    {
        lits = no-value;
    }

    return solved;
}
\end{lstlisting}

As you can see, all we've done is move our equation evaluation from brute force
into its own function. One thing to note is that this method also handles
updating our \textit{lits} field. Thus, if \textit{evaluate()} returns true, then
\textit{lits} will hold the satisfying assignment, and if it returns false, then
\textit{lits} will hold nothing.

Now our recursive \textit{solve()} method:

\begin{lstlisting}
bool solve(bool[] assignment)
{
    if (assignment.size == numVars)
    {
        return evaluate(assignment);
    }

    append(assignment, True);

    if (solve(assignment))
    {
        return true;
    }

    assignment.back() = False;

    if (solve(assignment))
    {
        return true;
    }

    truncate(assignment, 1);
    return false;
}
\end{lstlisting}

This might be slightly confusing at first, so let's go over it.

We first check our base case. If our assignment has all of its variables, we
evaluate it and return the result whether it be true or false. And you'll notice
that whenever we return true from a call of \textit{solve()}, we immediately
return true afterwards, so therefore, if we find a satisfiable assignment, we
will go up our call stack back to \textit{solveCNF()}, so no need to worry that
it will terminate if we find a satisfiable assignment.

However, if our assignment still needs some variables, we append true to our
assignment, then we do a recursive call on \textit{solve()}. If that recursive
call leads to a satisfiable assignment, we return true, otherwise, we swap out
our last value for false and try \textit{solve()} again. And if that ends up
working out, we return true, but if it doesn't then we remove the last value
and return false.

So how do we know that it will guarantee termination for an unsatisfiable problem?
Notice how when assigning the current variable true or false and it doesn't
work, we just remove the last one and go up the call stack by one? Well
suppose we're at the top of the call stack (we've assigned the first variable)
and the first variable is assigned false, and that didn't lead to a satisfiable
assignment. Then we remove that first variable from the assignment, and then
we return false. So our assignment array is empty, and we've returned to
\textit{solveCNF()}. Thus, \textit{solve()} will terminate eventually.

And how do we know that we'll try every single assignment?
Well for the very last variable, it's easy to see, if it being true doesn't satisfy
the equation, we assign it false, and if that doesn't work, we remove it and go
up the call stack. Then for the previous variable, if it was true, we assign it
false and go back to the last variable where it tries true and/or false again.
Then that doesn't work, we go up the call stack twice, and that continues until
we've found a satisfiable assignment or we've exhausted them all.\footnote{
This is a very informal proof, but what's most important is that we build up
our understanding of what's going on instead of dumping notation. Perhaps I
will do a formal proof in the future.}

Lastly, our \textit{solveCNF()} method:

\begin{lstlisting}
bool solveCNF()
{
    bool[1..numVars] assignment;

    return solve(assignment);
}
\end{lstlisting}

Because we're dealing with pseudocode, we don't have to do things such as
reserving or resizing the assignment array. And since \textit{evaluate()} takes
care of making sure our \textit{lits} field holds a valid assignment (if found)
and holds nothing if not, we don't need to do anything else in
\textit{solveCNF()}.

You may wonder though, why not pass in \textit{lits} to \textit{solve()} or
have no setup and make \textit{solveCNF()} purely recursive?
Those are very valid questions to ask.

This has to do with the differences between pseudocode and actual
implementation. In the actual implementation, some methods depend on the
state of \textit{lits}, such as whether or not it holds anything, and by passing
it into \textit{solve()}, we force it to hold something, even if there is no
satisfiable assignment. And it's true, that in \textit{solveCNF()} if we had no
setup, we could make it recursive on \textit{lits}, however, I chose to
implement it differently.

When you are reading this, please ask yourself these kinds of questions,
because oftentimes, the way you read my implementation is not the best way,
and even if it is, there is no reason as to why you can't think of different
implementations, try them out yourself, and compare. Perhaps, and I expect
this will happen often, you will find a better implementation than the one
written down here.

\subsection{Recursive Backtracking -- Results}

Running this code on the same 1000 instances (uf20-91), it took
\textbf{1,846,644ms}, which is roughly 30.78 minutes. This is actually a pretty big
improvement from our 35.32 minutes via brute forcing. Although this seems
like a big gain in speed, it should be noted that at this point the algorithm is
still just a fancier brute forcing algorithm. The next time I rewrote this algorithm,
it took \textbf{1,540,613ms}, which is about 25.68 minutes. So once again, that's about a 
~5 minute speedup just from rewriting it. Perhaps we should rewrite every piece of 
software in the world to make it run 5 minutes faster. 

Here's a question for you, why might assigning our variables true instead of
false first lead to an increase in speed? \textbf{Hint: Look at the number of negations 
present in the testing set.}

\begin{FVerbatim}
On all 1000 equations in uf20-91:

Recursive Backtracking:    1,540,613ms (~25.68 minutes).
Brute Force:               1,770,851ms (~29.51 minutes).
Short-Circuited (Basic):   49,874ms (~50 seconds).
Parallel:                  10,589ms (~10.5 seconds).
\end{FVerbatim}

\section{Unit Propagation}

Unit propagation is the next step in implementing DPLL, and it introduces the
idea of implication. Imagine a clause with two variables: $A$ and $Z$,
and we know that $A = False$. What can we say about $Z$? Well, if our equation is satisfiable, 
then the clause must be satisfiable, therefore $Z = True$ because if it didn't, 
that would be a contradiction.

So instead of assigning variables $B$, $C$, \ldots, $Y$, then assigning $Z$,
why don't we just do it now, and use that knowledge in other clauses
containing $Z$ or $\neg Z$?

That is essentially the idea behind unit propagation: the decisions we make
imply other assignments, and so we should force them now instead of waiting
until later.

However, this introduces a light hiccup into our implementation. You see, 
our initial invariant about assigning variables is that they are
assigned in numerical order instead of (possibly) out of order.

``Not a problem! We'll just make the assignment array have space for every
variable so we can index it anywhere,'' I hear you think. Unfortunately that interferes with
our base case because it relies on the size
of the assignment array. If it is initially at a fixed size of
\textit{numVars}, then our base case will always fire because our assignment
array always has a size of \textit{numVars}. So we would either need to
change our base case or change how we store learned variables.

What about a learned set? One that holds whether or not we've learned a
variable (positive or negative), and then as we traverse down the call stack,
we just check if a variable is already learned, and if so, we give it its learned
value?

That would 100\% work. In fact, the first time I implemented DPLL,
that is exactly what I used. However, it has the disadvantage of being annoying 
to maintain. Instead, why don't we create a new data type called something like
\textit{var} that represents the state of each variable: \textit{True},
\textit{False}, and \textit{Unassigned}.

Then we can just maintain a simple counter of how many are not
\textit{Unassigned}, and then that keeps our base case nice and neat.
Yes dear reader, we are going to do that, however, if you'd like to use the
learned set, reading about how we do it without a learned set will help you
and offer clues on how to implement it your way.

\subsection{Changes To Our Implementation}

\begin{lstlisting}
    public interface:
        SATSolver(string input)
        bool solveCNF()

    private interface:
        enum var
        {
            False
            True
            Unassigned
        }

        void parse(string equation)
        bool solve(var[] assignment, int level)
        bool evaluate(var[] assignment)
        void propagate()

        int numVars
        int numClauses
        int numAssigned

        int[][] clauses

        optional var[] vars
\end{lstlisting}

Here we just add in a new enum for the states that a variable can be, and
adjust the methods using \textit{bool}s to use the new enum type instead.

\subsection{``Textbook" Unit Propagation}

Before we start changing our algorithm, let's implement the
\textit{propagate()} method.

As stated before, the idea is that given a clause with either only one variable
or only one \textit{Unassigned} variable (and the rest evaluate to
\textit{False}), we can immediately say that \textit{Unassigned} variable will
be whatever value it needs to be to evaluate to \textit{True} for that clause.

In the ``textbook" way of implementation,
\footnote{\url{https://en.wikipedia.org/wiki/Unit_propagation}} 
once a clause is satisfied because of a literal, you remove every other
clause that is satisfied by that literal. And for every other clause containing 
the negation of the literal, you remove that negation entirely. This is done to shrink the 
search space.

We will implement this algorithm with as much malicious compliance as you can 
carry at an office job without getting fired:

\begin{lstlisting}
    bool solveCNF() {
        var[1..numVars] assignment

        for each (var in assignment) {
            var = Unassigned
        }

        propagate(assignment)
        solve(assignment, 0)
    }
\end{lstlisting}

We want to propagate before we start solving in case there are clauses with 
only one variable (it's already unit).

\begin{lstlisting}
    bool solve(bool[1..numVars] Assignment, int level) {
        if (level == numVars - 1) {
            return evaluate(Assignment)
        }

        bool[1..N] backupAssignment = assignment
        int[][] backupClauses = clauses

        if (Assignment[level] == Unassigned) {
            Assignment[level] = True
            propagate(Assignment)
        }

        if (solve(Assignment, level + 1)) {
            return True
        }

        clauses = backupClauses
        Assignment = backupAssignment

        if (Assignment[level] == Unassigned) {
            Assignment[level] = False
            propagate(Assignment)
        }

        if (solve(Assignment, level + 1)) {
            return True
        }

        return False
    }
\end{lstlisting}

Before we get onto the code for unit propagation, let's understand this code first. 
Because unit propagation no longer guarantees that variables are assigned in order, 
we need to pass in a ``level'' as a parameter so we know which variable we need to 
assign. We also now need to check that propagation didn't already assign this variable. 
If it's currently Unassigned, we need to assign it then propagate, but if it's been 
assigned before we assign it, then we know that it's the result of propagation, so we 
can continue onto the next level. Other than that, it's practically the same as our recursive 
backtracking code, except now we need to have a ``backup'' of our clauses and our assignment. 

Why? Because \textit{propagate()} \textbf{mutates} the clauses and the assignment in a way 
that breaks the assumption that our literals are assigned in ``numeric'' order. Since it 
returns nothing, we don't have a way to undo that easily, so instead, for readability, 
we just have a backup copy and whenever propagation leads to an unsatisfying assignment, 
we return false and undo whatever changes we made by copying back into our solver's fields.
And the reason why we don't need to undo the propagation before returning False is because 
when we move back up the call stack by one level, the method continues by undoing the effects for us.
Now let's look at how we perform textbook unit propagation:

\begin{lstlisting}
    void propagate(int[1..N] Assignment) {
        while (true) {
            Set<int> unitClauses;
            Set<int> unitVars;

            for (i = 0 up to numClauses) {
                clause = clauses[i]

                bool isUnit = True
                int unassignedIdx = -1
                int numUnassigned = 0

                for (j = 0 up to clause.size) {
                    int var = clause[j]
                    int idx = abs(var)

                    if (Assignment[idx] == Unassigned) {
                        numUnassigned = numUnassigned + 1
                        unassignedIdx = j
                    } else if (var > 0 && clause[j] == True 
                        || var < 0 && clause[j] == False)
                    {
                        isUnit = false
                    }
                }

                isUnit = isUnit && numUnassigned == 1

                if (isUnit) {
                    add(unitClauses, i)
                    add(unitVars, clause[unassignedIdx])

                    int var = clause[unassignedIdx]
                    int idx = abs(var)

                    if (var > 0) {
                        assignment[idx] = True
                    } else {
                        assignment[idx] = False
                    }
                }
            }

            int clauseIdx = 0;
            for each (clause in clauses) {
                for each (var in clause) {
                    if (contains(unitVar, -var)) {
                        remove(clause, -var)
                    }

                    if (contains(unitVar, var) 
                        && contains(unitClause, clauseIdx))
                    {
                        remove(clauses, clause)
                    }
                }

                clauseIdx = clauseIdx + 1
            }

            if (unitClauses.size == 0 && unitVars.size == 0) {
                break
            }
        }
    }
\end{lstlisting}

Okay, now let's understand what's happening here. We immediately go into an 
infinite loop (nice), and we start looking at each clause for two things:
\begin{enumerate}
    \item The clause has only one Unassigned variable 
    \item Every other literal in the clause is False (remember, literals are evaluated variables)
\end{enumerate}

If both things are true, then the Unassigned variable will be assigned whatever value 
is required to make it evaluate to True within that clause, and we know we have found 
a unit variable and a unit clause. After we have found all of our unit variables and unit 
clauses, we then rescan every other clause and delete the each unit variable's negative,
(if we learn that -5 is unit, then we remove +5), and if the clause contains a unit 
variable, we just delete that clause entirely from the database. Once we stop finding 
new unit variables and new unit clauses, we stop entirely. You may wonder (I reread my 
code and thought I found a bug), ``how do we know that we always terminate?'' That's a 
great question. You see, this confusion may come from the fact that we rescan the database 
every iteration of the \textit{while()} loop searching for unit clauses. What is stopping 
us from re-finding the same clause as unit upon a new iteration? The answer is much simpler 
than you'd expect: since we assign our unit variables, the clauses upon the next iteration 
no longer foloow our criteria for being unit propagate-able because there is now a literal 
in that clause that is True. Therefore, the number of unit variables and clauses monotonically
decrease between iterations (a fancy way of saying the number can only go down).

Now, you may be wondering, ``what's the performance of our solver now?" Let me answer that 
as plainly as possible: worse than brute force. How? Easy! We are doing SO MANY things that 
hurt the computer. We are allocating TONS of memory in order to backtrack, we are allocating 
even more memory to propagate (hash sets aren't cheap), and hashing, while O(1), is not 
a cheap operation to do. So yeah, I can believe that it's worse than brute force because 
at least the brute force algorithm was easy on the CPU. Let's learn how this is 
\textit{usually}
\footnote{I wanted to say ``actually implemented'', however, that's a big generalization, and 
again, since I'm not the expert here, I don't think I have the authority to say that. The next 
section is how \textit{I} implemented it, and it roughly follows the idea on how SAT solvers 
implement this.}
implemented!

\section{Better Unit Propagation}

So the textbook style unit propagation sucked. Let's change that. As discussed before, 
the biggest bottleneck/inefficiency in the actual implementation is the number of copies 
we need to make of our clauses and assignments on each decision level. Assigning one variable 
creates a new decision level, so 20 variables means 20 levels. For an unsatisfiable 
CNF equation, I estimate we need to make about 3 copies of our clauses and assignments:
one for the initial storage, another for the first failure, and a third for the second. 
Even if we ignore the inefficiencies in how we propagate (memory allocations, traversals, and 
adjustments our arrays would need to make when we delete things), this is a huge amount of 
memory to store. If you think that 20 variables and 91 clauses is a lot, modern SAT solvers 
can handle equations with \textbf{millions} of variables, so creating an algorithm that 
makes brute force look attractive is actually an achievement.
\footnote{They might revoke my degree for that actually}

So what's the strategy? Firstly, it would be nice if we didn't actually mutate our 
clauses at all. Ie, no clause deletion, and no variable deletion. Next, it would be 
nice if we didn't need hash sets in order to store our unit variables. Instead, it
would be nice if we had a data structure where we put in our decided variables, and then 
it slowly shrinks as we iterate through. Ie, we propagate one on variable at a time and 
never again. This sounds like we need a \textit{queue}.
\footnote{That wasn't shoehorned in at all!}

Lastly, instead of copying our assignment, and constantly reiterating through the 
\textit{entire} clause database, what if we just marked which variables we assigned,
and which clauses we've satisfied? That way when we need to undo, we just mark those 
variables as Unassigned, and those clauses as no longer satisfied. That way we save a lot 
of memory (and besides, copying is NOT O(1) as you might be led to believe by that deceptively
simply ``='' operator).

Let's change our interface:

\begin{lstlisting}
    public interface:
        SATSolver(string input)
        bool solveCNF()

    private interface:
        enum var
        {
            False
            True
            Unassigned
        }

        void parse(string equation)
        bool solve(var[] assignment, int level)
        bool evaluate(var[] assignment)
        (int[] vars, int[] clauses) propagate()

        int numVars
        int numClauses
        int numAssigned

        int[][] clauses
        bool[] satisfied

        optional var[] vars
\end{lstlisting}

We need to make some minimal changes to \textit{solveCNF()}:

\begin{lstlisting}
    bool solveCNF() {
        var[1..numVars] assignment

        for each (var in assignment) {
            var = Unassigned
        }

        propagate(assignment, 0)
        solve(assignment, 0)
    }
\end{lstlisting}

Since \textit{propagate()} now takes a decision level/variable index, we need to pass 
one to the method. Although no variable has an index of 0 (because it's reserved to mark the 
end of a clause), passing in $0$ just tells it ``Hey, there isn't a literal to propagate on,
so just find all clauses with only one variables and mark them as satisfied and `learn' those 
variables.'' And since this is at the root level, we don't need to undo anything if it turns out 
that the equation is unsatisfiable, since we're gathering things that HAVE to be a certain value 
if this equation is to be satisfiable.

Since we don't need to change \textit{solveCNF()}, we can go straight to \textit{solve()}:

\begin{lstlisting}
    bool solve(int[1..N] Assignment, int level) {
        if (level == numVars - 1) {
            return evaluate(Assignment);
        }

        if (Assignment[level] != Unassigned) {
            return solve(Assignment, level + 1)
        }

        Assignment[level] = True
        (int[] unitVars, int[] satisfiedClauses) 
            = propagate(Assignment, level)

        if (solve(Assignment, level + 1) == True) {
            return True
        }

        for each (var in unitVars) {
            int idx = abs(var)
            Assignment[idx] = Unassigned
        }

        for each (clauseIdx in satisfiedClauses) {
            satisfied[clauseIdx] = False
        }

        Assignment[level] = False
        (int new_unitVars, int new_satisfiedClauses) 
            = propagate(Assignment, level)

         if (solve(Assignment, level + 1) == True) {
            return True
        }

        for each (var in new_unitVars) {
            int idx = abs(var)
            Assignment[idx] = Unassigned
        }

        for each (clauseIdx in new_satisfiedClauses) {
            satisfied[clauseIdx] = False
        }

        Assignment[level] = Unassigned
        return False
    }
\end{lstlisting}

This method is practically the same, except now, we just check if we've already assigned 
the variable at the current decision level. If we have assigned a variable before we reach 
its decision level, then it must have been inferred/learned from the \textit{propagate()} 
method, so we don't need to call it again. Then, it's business as usual: assume the 
current variable is True, propagate, if it leads to a satisfiable assignment, return True, 
otherwise, we undo what we've assigned and marked as learned, then repeat with False. 

What's more interesting is our shiny new \textit{propagate()} method:

\begin{lstlisting}
    (int[] vars, int[] clauses) propagate(int[] assignment, int level) {
        queue<int> unit

        int[] unitVars
        int[] satisfiedClauses 

        unit.enqueue(level)

        while (unit.size > 0) {
            unit.dequeue()

            for (i = 0 up to numClauses) {
                if (satisfied[i]) {
                    continue
                }

                clause = clauses[i]

                bool isUnit = True
                int unassignedIdx = -1
                int numUnassigned = 0

                for (j = 0 up to clause.size) {
                    int var = clause[j]
                    int idx = abs(var)

                    if (Assignment[idx] == Unassigned) {
                        numUnassigned = numUnassigned + 1
                        unassignedIdx = j
                    } else if (var > 0 && clause[j] == True 
                        || var < 0 && clause[j] == False)
                    {
                        isUnit = false

                        satisfied[i] = true
                        append(satisfiedClauses, i)
                    }
                }

                isUnit = isUnit && numUnassigned == 1

                if (isUnit) {
                    add(unitClauses, i)
                    add(unitVars, clause[unassignedIdx])

                    int var = clause[unassignedIdx]
                    int idx = abs(var)

                    if (var > 0) {
                        assignment[idx] = True
                    } else {
                        assignment[idx] = False
                    }

                    unit.enqueue(var)

                    append(unitVars, var)
                    satisfied[i] = true
                    append(satisfiedClauses, i)
                }
            }
        }

        return (unitVars, satisfiedClauses)
    }
\end{lstlisting}

So, let's discuss what's going on. \textit{propagate()} now takes in the decision level, 
which is also the same thing as an index for a variable (conceptually speaking). It then 
adds that to a queue, then while the queue is not empty, it scans the clauses as before. 
First it makes sure we're not operating on a clause that's already been satisfied 
(because we can't tell anything from it), and when it finds a clause that has a true literal, 
it marks the clause as satisfied. Then when it finds a unit clause, it marks the clause as 
satisfied, keeps track of the new unit variable its found, and then adds that unit variable 
to the queue. Then once it's done, it returns the unit variables and the indices of the satisfied 
clauses.

You may have some questions about this approach. Firstly, how do we make sure we don't 
ever add the same variable twice? The answer is the same as before: because we assign a 
variable once it's found to be in a unit clause, it can't be found as Unassigned again, 
and therefore can't be marked as a unit variable again. Then, since we skip over already 
satisfied clauses, the amount of clauses we scan each iteration decreases monotonically. 

And going back to our \textit{solve()} code, even if we mark the same clause as satisfied 
or the same variable as unit twice, it doesn't matter because when we backtrack, we simply 
restore it back to its state of unsatisfied or Unassigned, respectively. What's important is 
that we don't mark the same clause or variable twice \textit{across decision levels} because 
then, and only then, would our backtracking code become a problem because we wouldn't be 
correctly restoring the state of the equation back to what it was before the decision level 
was entered. But since we can't mark the same variable or clause twice \textbf{until}
we backtrack, there's nothing to worry about.

Before we end this section, you may also be wondering: can we use the \textit{satisfied}
array in our \textit{evaluate()} method to skip clauses we know we've already satisfied?
And the answer is yes, and I would encourage you to do so in your implementations.

\subsection{Better Unit Propagation - Results}
So after implementing better unit propagation, this took \textbf{868ms} to solve all 1000 
satisfiable equations. The first time I implemented this version of unit propagation,
it actually took \textbf{2363ms}, so this is somehow over twice as fast. And because it's so 
fast, just like the first time I implemented this, we will be now testing our solver on 
unsatisfiable equations as well.

Our second testing set comes from the same database as before.
\footnote{\url{https://www.cs.ubc.ca/~hoos/SATLIB/benchm.html}, in case you needed a reminder.}
And our next testing set is called ``UUF50.218.1000." In other words, it's full of equations that 
have 50 variables, 218 clauses, and there are 1000 such equations in the set. In order to solve 
all 1000 equations, it took \textbf{140,411ms} (~2.34 minutes). That's also much faster than the first time 
I implemented this algorithm. The first time I did, it took 646,228ms (~10.77 minutes). Honestly,
with this much speedup, I don't know if I was just bad at implementing these algorithms the first time,
\footnote{Highly likely}
or if downgrading my operating system (and subsequently wiping my harddrive clean)
\footnote{I HATED MacOS Tahoe. Sequoia is where it's at.}
just has that much of an effect on execution speed.
\footnote{Also likely, depending on how much RAM and swap is being used to begin with.}

Why didn't we test our previous versions out on unsatisfiable equations? That's a good question.
For brute forcers, since there is no satisfiable assignment, it needs to check every single 
one to conclude that it's unsatisfiable, which would take a lot of time. Likewise with our 
recursive backtracking solver. And since our (my) implementation of textbook unit propagation 
was \textit{worse} than brute force,
\footnote{A difficult feat.}
obviously we wouldn't want to test it on something that will take more time than a satisfiable 
equation. Also because the asymptotic runtime is \textit{still} exponential, having a much 
larger equation would take a much, much larger amount of time, even if they were satisfiable.

\begin{FVerbatim}
Results so far:

On all 1000 equations in uf20-91:

Unit Propagation:          868ms (~0.87 seconds).
Recursive Backtracking:    1,540,613ms (~25.68 minutes).
Brute Force:               1,770,851ms (~29.51 minutes).
Short-Circuited (Basic):   49,874ms (~50 seconds).
Parallel:                  10,589ms (~10.5 seconds).

On all 1000 equations in UUF50.218.1000:

Unit Propagation: 140,041ms (~2.34 minutes).
\end{FVerbatim}

\section{Pure Literal Elimination}

Suppose I gave you the following CNF equation:

\begin{FVerbatim}
1 2 -3 0
1 -4 5 0
1 -2 5 0
\end{FVerbatim}

Do you notice anything odd about it? Perhaps ``odd" is not the correct word here, do you notice 
an easy way to satisfy this equation? You probably noticed that the variable $1$ shows up in every 
clause and is always in the positive. So, if you just assign $1 = T$, then it doesn't matter what 
the rest of the variables are assigned, because $1 = T$ satisfies every clause! This is idea behind 
\textbf{Pure Literal Elimination}. Essentially, every now and then, we scan our equation for ``pure 
literals'', which are variables that only show up in the positive or negative (but never both), and
then give them whatever assignment is necessary to have their literals evaluate to True. 

Now, suppose I give you this equation:
\begin{FVerbatim}
1 0
-1 -2 -3 0
1 2 0
-2 3 0
\end{FVerbatim}

I have a question for you dear reader: is $-2$ a pure literal? You might be inclined to say ``no"
because $-2$ and $2$ are both present in the equation. However, despite that, the answer is ``yes.''
Why? Well notice in the first clause that $1$ is a unit variable, so $1 = T$ is forced by pure literal 
elimination. Then the third clause, $1 2 0$, because it's already satisfied, the variable $2$
doesn't matter anymore. The other two clauses, which are unsatisfied, are the ones that matter. 
They both contain $-2$, and because they're unsatisfied, we can force $2 = F$ to satisfy them,
thus making the assignment of $3$ irrelevant. 

Okay, that makes sense, sort of. How come we're allowed to do that? Like, how do we know that 
forcing $2 = F$ wouldn't screw up the equation elsewhere (if it were larger)? And the answer is:
because it doesn't hurt the other clauses. In the third clause, once it's satisfied, the only thing 
that can make it unsatisfied is the backtracking, but since $1$ is unit, there shouldn't be any 
backtracking that would suddenly make its satisfiability rest on the assignment of $2$. Likewise, 
for the other two clauses, no matter what value $3$ takes, $2 = F$ must be true to satisfy the 
other clause (try it out for yourself!), so we can go straight ahead and assign $2 = F$ and skip 
deciding. 

Okay, sure, but there are 8 total assignments that we can get. How do we know that forcing 
$2 = F$ would lead to a satisfiable assignment? As in, how do we know we're not prematurely 
boxing ourselves in? That's another great question! The reason why we're not boxing ourselves 
in prematurely is because, again, it only helps the satisfiability of the whole equation, 
and because we're not doing anything detrimental, we can make that decision.
\footnote{Essentially, think of it as making the solution ``more correct'' 
by forcing the pure literal.}

\subsection{Implementing Pure Literal Elimination}

Now that we (sort of) understand why it works, how do we implement it?
We need some form of efficient tracking for things we've already encountered, perhaps 
something like a map or a set. Then, we just need to make sure we've only encountered 
one version of a literal, and if we have, then we know it's pure. Then we just need 
to gather our pure literals and the satisfied clauses, and then return them. 

\begin{lstlisting}
(int[] var, int[] clauses) pureElim(int[] assignment) {
    map<int, int[]> var_to_clause
    int[] pure
    int[] satisfied

    for (i = 0 up to numClauses) {
        if (satisfied[i]) {
            continue
        }

        clause = clause[i]
        for each (var in clause) {
            append(var_to_clause[var], i)
        }
    }
    
    for each ([var, clauses] in var_to_clause) {
        idx = abs(var)

        if (contains(var_to_clause, -var) or assignment[idx] != Unassigned) {
            continue
        }

        if (var > 0) {
            assignment[idx] = True
        } else {
            assignment[idx] = False
        }

        for (i = 0 up to clauses.size) {
            satisfied[clauses[i]] = true
            append(satisfied, clauses[i])
        }

        append(pure, var)
    }

    return (pure, satisfied)
}
\end{lstlisting}

So this is a single pass version of pure literal elimination. It is possible that 
after finding some pure literals, and marking some clauses satisfied, a second pass could 
find new pure literals. So why are we implementing a single pass version? Besides 
it not being a technique used in modern SAT solvers, in my experience, which isn't much 
since I'm being honest, it doesn't make a big (positive) difference in solve time 
efficiency. And besides, if we wanted a multi-pass version, we could wrap it in 
\textit{while(true)} loop and break once we find no more new pure literals. 

Now let's get into how the algorithm works. We first build a map from variables 
to their clauses (in unsatisfied clauses). Next, afterwards, we then check 
for each element in the map, that the map does not contain the negative variable 
of the current entry, and if it doesn't, we then need to make sure that the current
variable is Unassigned. If the current entry passes both those conditions (pure and 
Unassigned), we assign it so it evaluates to True. Then, we mark each clause the 
variable is found in as satisfied, and then we add the indices of the satisfied clauses 
into the return values. 

\subsection{Pure Literal Elimination - Results}

So. I have some terrible news. After implementing pure literal elimination, our solving 
time \textit{\textbf{increased}}. I know, I know. It's not supposed to happen, but it can 
happen. There are many possible explanations for this. For instance, pure literal elimination
takes a decent amount of work to perform. We have to scan all unsatisfied clauses, and since 
there is no way to determine if a literal is impure until after everything is scanned
(we could have used a hash set, but then we'd need to hash for every literal), we're 
accruing memory and runtime costs. Second, even if pure literal elimination finds something,
it usually isn't much (at least for the testing sets currently in use), so we're doing 
so much work for so little in return. Third, it increases our backtracking cost because now 
we need to undo unit propagation and the pure literal elimination. Fourth, my implementation
isn't perfect. It is possible that forcing a pure literal may force another pure literal 
in a different clause, so you would need multiple passes to do. However, our implementation 
is just one pass (wrap it in a \textit{while} loop or something).
\footnote{Modern SAT solvers don't tend to use pure literal elimination for reasons outside 
the ``it wasn't very helpful'' aspect, and next chapter we'll see why.}

However, no one cares \textit{why} something performed badly. People only care about 
\textit{how} badly something performed. So here are the results:
For uf20.91: 1548ms, almost a 2x increase from just unit propagation.
For UUF50.218.1000: 231,107ms (~3.85 minutes), about a 1.5 minute increase 
from just unit propagation.

\begin{FVerbatim}
On all 1000 equations in uf20-91:

Pure Literal Elimination:  1548ms (~1.55 seconds).
Unit Propagation:          868ms (~0.87 seconds).
Recursive Backtracking:    1,540,613ms (~25.68 minutes).
Brute Force:               1,770,851ms (~29.51 minutes).
Short-Circuited (Basic):   49,874ms (~50 seconds).
Parallel:                  10,589ms (~10.5 seconds).

On all 1000 equations in UUF50.218.1000:

Pure Literal Elimination: 231,107ms (~3.85 minutes).
Unit Propagation:         140,041ms (~2.34 minutes).
\end{FVerbatim}

Pure Literal Elimination had about a \textbf{2x} increase (which is bad) in solving time
when compared to just Unit Propagation on uf20.91, which is the test set of satisfiable
problems. It also had about a \textbf{1.6x} increase in solving time on UUF50.218.1000,
which is the set of unsatisfiable problems.

\section{Watched Literals}

Before we get into the Watched Literal algorithm, I want to ease us into it.
Why? Because it's so different from our current version of unit propgation. 
For a moment (and for the rest of the book) let's forget about pure literal 
elimination; it wasn't that great anyways. Unit propagation gave us a massive speedup 
oppossed to recursive backtracking, taking runtimes down from minutes to seconds and 
from half-hours to mere minutes. Despite all of that, it's really inefficient in the 
sense that we're going so much work for so little in return.

Let's define what a \textbf{contradiction} is, because we've been using that word 
here and there, but beyond ``unsatisfiable,'' we haven't formally defined what it means.
A \textbf{contradiction} is the result of a literal forcing a clause to evaluate to False,
leading to the whole assignment to become unsatisfiable. You might remember how in the 
brute force solver, we added an \textit{if()} statement at the end of evaluating a singular 
clause that checked if that clause evaluated to False. If it did, we skipped that assignment 
and moved onto the next one, but if it didn't, we continue our evaluation. 

This led to a pretty significant speedup in the brute forcer because instead of evaluating everything 
then deciding, we did it incrementally, allowing us to quickly leave when something was bad. 
At the moment, our original idea for unit propagation follows that first idea of delaying 
evaluation to the end. However, if you go back and look at the pseudocode, you'll notice that we 
do something \textit{almost} the same as checking if a clause evaluates to False, at the moment,
we check if everything except for one element evaluates to False, and that one element is Unassigned
(since we made it so that unit propagation only operates on unsatisfied clauses). However, notice 
that when a clause evaluates to False, rather than stopping our propagation and assignment, we 
continue until we reach the evaluation phase? 

Just like in brute force, the idea is the same: once a clause evaluates to False, no matter 
where we are in our assignment, there's no point in continuing on because we've found out that 
nothing will unfalsify that clause. As an exercise, I want you to try adding that into your 
unit propagation code and seeing how fast it takes to run on our test sets. I would do it, but 
I have a \underline{massive} headache from implementing this algorithm last night.
\footnote{This headache was also induced by my lack of ability to fall asleep at all last night
AND a really intense session at the gym. Computer science gives you headaches for sure, but not like 
this one.}
Once you're done, we'll discuss watched literals some more.

So that early contradiction detection is one of the ideas behind watched literals. 
It's not one the key ideas, but it is part of it. There are some other ideas behind 
watched literals. For instance: zero cost backtracking. Not O(1) backtracking, $0$
backtracking. Speaking of cost, do you notice something else wasteful about the 
way we propagate? If not for our \textit{satisfied} marker array, we'd have to 
rescan the whole equation every time we propagated, which is about twice per level. 
That's pretty wasteful. After all, why do have to scan every clause in their entirety
every single time? Especially because, usually, clauses don't become unsatisfied very 
often. And because they're not at risk of becoming unsatisfied very often, why do we need 
to scan it? And besides, it's not like every single variable appears in every single 
clause, so why not scan only the clauses we have to? I hope these questions are leading 
you to think of other questions regarding inefficiencies in our current implementation.

These questions were the driving force into the invention of the 
\textit{Watched Literals} algorithm.

\textbf{History can be found here}
\footnote{\url{https://school.a4cp.org/summer2011/slides/Gent/SATCP3.pdf}}

Let's learn how the Watched Literals algorithm works. Earlier, I made a comment 
about how usually, each clause doesn't contain every single variable in the 
equation. So when we force a variable's assignment, we don't need to check 
every single clause, only at most the ones that contain the variable and its 
negation. (Ie, if we force a variable, say $2 = T$, we only need to check the clauses 
containing $2$ and the clauses containing $-2$). However, do we \textit{need} to 
check the clauses containing $2$? Think about it. Those clauses are already 
satisfied by $2 = T$, so there's nothing to propagate on, for those clauses. 
The only clauses that have a \textit{chance} to produce a propagation are the 
ones containing $-2$ because $-2 = F$, so if they have one Unassigned variable 
and everything else is now False, we can propagate!

\subsection{A Visual Example of Watched Literals}
Let's view a visual example. Suppose we had a clause with 5 literals 
(doesn't matter what they are). The first two literals are watched
\begin{center}
    \[ \blackinwhitesquare \blackinwhitesquare \square \square \square \]
    \[ U U U U U \]
\end{center}
Suppose the first literal now evaluates to False. We now scan the clause for 
another literal that doesn't evaluate to False. We will find it at the third literal.
\begin{center}
    \[ \square \blackinwhitesquare \blackinwhitesquare \square \square \]
    \[ F U U U U \]
\end{center}
Now suppose the fifth literal becomes False. We do nothing because we're not watching 
the fifth literal.
\begin{center}
    \[ \square \blackinwhitesquare \blackinwhitesquare \square \square \]
    \[ F U U U F \]
\end{center}
Now the third literal becomes False. We scan again for another not-False literal.
\begin{center}
    \[ \square \blackinwhitesquare \square \blackinwhitesquare \square \]
    \[ F U F U F \]
\end{center}
Now the second literal becomes False. Other than the second watched literal (the fourth),
there are no more not-False literals. We don't relocate the watch.
\begin{center}
    \[ \square \blackinwhitesquare \square \blackinwhitesquare \square \]
    \[ F F F U F \]
\end{center}
However, now the second watched literal knows it's the last Unassigned literal in the clause.
It will now evaluate to True!
\begin{center}
    \[ \square \blackinwhitesquare \square \blackinwhitesquare \square \]
    \[ F F F T F \]
\end{center}
Our clause is now satisfied. Suppose instead, when we were here:
\begin{center}
    \[ \square \blackinwhitesquare \square \blackinwhitesquare \square \]
    \[ F F F U F \]
\end{center}
That the fourth literal also became False.
\begin{center}
    \[ \square \blackinwhitesquare \square \blackinwhitesquare \square \]
    \[ F F F F F \]
\end{center}
We've now encountered a \textit{contradiction}, meaning that whatever assignment we have
is impossible to satisfy the equation. Suppose we backtrack, and that restores us to a 
state where two literals in this same clause are now Unassigned. What happens to the 
watched literals?
\begin{center}
    \[ \square \square \square \square \square \]
    \[ U U F F F \]
\end{center}
The answer is that they stay where they are.
\begin{center}
    \[ \square \blackinwhitesquare \square \blackinwhitesquare \square \]
    \[ U U F F F \]
\end{center}
This is because one watched literal is Unassigned, so everything is okay. Although 
this last diagram isn't accurate to what a real restored state would look like, we can
suspend our disbelief for the sake of simplicitly. The reason why watched literals work is 
because of the \textbf{Watched Literal Invariant}, which is that at any given point in time,
at least one of the watched literals is either True or Unassigned. This invariant is what
allowed us to propagate the clause when there was only one Unassigned literal left, and we 
will see further how it helps.


\subsection{Algorithmic Changes}
\begin{lstlisting}
public interface:
    SATSolver(string input)
    bool solveCNF()

private interface:
    enum var
    {
        False
        True
        Unassigned
    }

    void parse(string equation)
    bool solve(var[] assignment, int level)
    optional int[] delta propagate(var[] assignment)
    bool createWatched()

    int numVars
    int numClauses

    int[][] clauses
    bool[] satisfied
    queue<int> unitProp

    optional var[] vars
    map<int, int[]> var_to_clause
    map<int, (int, int)> clause_to_vars
\end{lstlisting}

So firstly, there are some major differences. We removed the \textit{evaluate()} method. 
This is because now that propagation checks for contradictions, having all variables 
assigned is now \textbf{\textit{equivalent}} to finding a satisfactory assignment of 
variables. So we don't need to evaluate our assignment anymore! That's significantly 
less work than what we had before, which is what we want. Next, we added in a queue, 
\textit{unitProp}, kind of like what we had for unit propagation before. This time,
it's now a member of the solver object instead of a local object within the method.
We will get back to that in a second. Lastly, we added in two maps, one from 
variables to clauses, and the other from clauses to a pair of variables. 

These data structures are used for fast lookup between variables to  
\textbf{\textit{indices}} of clauses (for \textit{var\_to\_clause}), and 
one from \textbf{\textit{indices}} of clauses to \textbf{\textit{indices}} of 
watched literals. It's important that we remember we're using indices instead 
of actual clauses because forgetting how we're using our mapping is going to create 
a lot of bugs and we're gonna have a bad time. 

We also added in a new method, called \textit{createWatched()} and what it does 
is it creates our initial maps for var\_to\_clause and clause\_to\_vars. 
Here is the pseudocode:

\begin{lstlisting}
    bool createWatched(int[] assignment) {
        for (i = 0 up to numClauses) {
            clause = clauses[i]

            if (clause.size >= 2) {
                append(var_to_clause[clause[0]], i)
                append(var_to_clause[clause[1]], i)

                clause_to_var[i] = {0, 1}
            } else if (clause.size == 1) {
                int var = clause[0]
                int idx = abs(var)
                append(var_to_clause[var], i)
                clause_to_var[i] = {0, 0}

                enqueue(unitProp, var)

                if (var > 0) {
                    assignment[var] = True
                } else {
                    assignment[var] = False
                }
            } else {
                return False
            }
        }

        return True
    }
\end{lstlisting}

\noindent So here's how to read it: for each clause, it will fall into one of three cases:
\begin{enumerate}
    \item It's size is greater than or equal to 2
    \item It only has 1 variable
    \item It's empty
\end{enumerate}

In the first case, the two literals that will be watching that clause will be the 
first two literals. So the index $i$ gets appended to their array of clause indices.
Then, we record that clause $i$ is being watched by its first and second literal.

In the second case, since there's only one literal, it gets watched by only that 
literal (duh). We also record that clause $i$ is being watched by only its first literal.
And lastly, we enqueue that literal to our unit propagation queue so we can do a root 
level unit propagation before we start solving. We also want to immediately force the 
assignment of the variable so that way we can use its value immediately in solving.

In the last case, since there are no literals, this clause is automatically unsatisfiable, 
so therefore this whole equation is unsatisfiable, so we return False so that when we call 
this method, we know that we can't solve it. 

Now let's look at our modified solving methods:

\begin{lstlisting}
    bool solveCNF() {
        var[1..numVars] assignment

        for each (var in assignment) {
            var = Unassigned
        }

        if (createWatched(assignment) == False) {
            return False
        }

        if (propagate(assignment, 0) returns nothing) {
            return False
        }

        solve(assignment, 0)
    }
\end{lstlisting}

There are some slight differences: we create our watched literals and make sure 
that there isn't an unsatisfiable clause, and we make sure that \textit{propagate()}
returns something. If both those things are successful, then we can start solving away.

\begin{lstlisting}
    bool solve(int[] assignment, int level) {
        if (level == numVars) {
            return True
        }

        if (assignment[level] != Unassigned) {
            return solve(assignment, level + 1)
        }

        int cur_var = level + 1

        assignment[level] = True 

        optional int[] prop_result = propagate(assignment, cur_var)

        if (prop_result is optional) {
            assignment[level] = False

            optional int[] new_prop_result 
                = propagate(assignment, -cur_var)

            if (new_prop_result is optional) {
                return False
            }

            if (solve(assignment, level + 1)) {
                return True
            }

            int[] new_delta = new_prop_result.value

            for each (int var_idx in new_delta) {
                assignment[var_idx] = Unassigned
            }

            return False
        }

        if (solve(assignment, level + 1)) {
            return True
        }

        int[] delta = prop_result.value

        for each (int var_idx in delta) {
            assignment[var_idx] = Unassigned
        }

        assignment[level] = False

        optional int[] new_prop_result 
            = propagate(assignment, -cur_var)

        if (new_prop_result is optional) {
            return False
        }

        if (solve(assignment, level + 1)) {
            return True
        }

        int[] new_delta = new_prop_result.value

        for each (int var_idx in new_delta) {
            assignment[var_idx] = Unassigned
        }

        return False
    }
\end{lstlisting}

This is almost the same as our old unit propagation code, except for that we don't 
keep track of which clauses are satisfied across levels, and the propagation method 
returns an array of variable indices instead of the actual variables. The code 
is also split up into multiple branches depending on whether or not propagation was 
successful. Some parts of the code seem really similar to each other because of this. 
I am sure that there may be a way to simplify the similar parts, but I like the way I 
wrote it because I find it readable and the first time I wrote it, it was very nested.

Now for the interesting part:
\begin{lstlisting}
    optional int[] propagate(int[] assignment, int decision) {
        int[] delta

        if (decision != 0) {
            enqueue(unitProp, decision)

            int idx = abs(decision)
            append(delta, idx)
        }

        while (unitProp.size != 0) {
            int prop = unitProp.pop()

            if (contains(var_to_clause, -prop) == False) {
                continue
            }

            int[] watched_clauses = var_to_clause[-prop]

            for (i = 0 up to watched_clauses.size) {
                clause_idx = watched_clauses[i]

                int[] clause = clauses[clause_idx]
                (int, int) watches = clause_to_var[clause_idx]

                if (clause[watches.first] != -prop) {
                    swap(watches.first, watches.second)
                }

                bool relocated = False
                for (j = 0 up to clause.size) {
                    if (j == watches.first or j == watches.second) {
                        continue
                    }

                    int unwatched = clause[j]
                    int unwatched_idx = abs(unwatched)

                    bool evals_false = 
                        (unwatched > 0 and assignment[unwatched_idx]
                            == False)
                        or
                        (unwatched < 0 and assignment[unwatched_idx] 
                            == True)
                    if (!evals_false) {
                        append(var_to_clause[unwatched], clause_idx)
                        erase_at(watched_clauses, i)
                        watches.first = j
                        relocated = True
                        break
                    }
                }

                if (relocated == False) {
                    i = i + 1

                    int other_watch = clause[watches.second];
                    int idx = abs(other_watch)

                    if (assignment[idx] == Unassigned) {
                        if (other_watch > 0) {
                            assignment[idx] = True
                        } else {
                            assignment[idx] = False
                        }

                        append(delta, idx)
                        enqueue(unitProp, other_watch)
                    } else {
                        bool evals_false =
                            (other_watch > 0 and assignment[idx] 
                                == False) 
                            or
                            (other_watch < 0 and assignment[idx] 
                                == True)

                        if (evals_false == False) {
                            continue
                        }

                        for (k = 0 up to delta.size) {
                            assignment[delta[k]] = Unassigned
                        }

                        clear(unitProp)

                        return null
                    }
                }
            }
        }
        return delta
    }
\end{lstlisting}

This is a huge method, so let's walkthrough it so we understand every aspect of it. 
If we look back at our \textit{solve()} code, we pass in \textit{cur\_var} and 
-\textit{cur\_var} as the \textit{decision} parameter into the propagation method. 
And in the beginning of \textit{solveCNF()}, we pass in $0$, to signify propagation at the 
root level. When that happens, we don't add anything into the queue, instead we just propagate 
on the literals that were forced during the creation of the watched literals. 

Like I said before, we only need to check the clauses containing the negative version of the 
literal being propagated on because those are the clauses that the actual unit propagation could 
be forced upon. However, we only need to actually check the clauses being watched by the negative 
version of the literal being propagated on. 

Next, we grab the indices of those clauses being watched and then for each one, we try 
to relocate the index of the first watched literal to a new literal that doesn't evaluate 
to False. Why only the first? Because swapping the two indices and operating on only the 
first index is far easier than creating a second branch dedicated to the second index.
If a literal doesn't evaluate to False, we first make that unwatched literal now watch 
the current clause, we then make the current watched literal stop watching the current clause,
and then we update the first index of the watched literals to be the index of the newly 
watching literal. That means we have successfully relocated the watched literal and no 
unit propagation needs to occur. Because we're modifying the array of clause indices, 
we don't want to increment the index variable in our implementation.

\textit{Aside:
When you delete something in a dynamic array, in order to fill the blank space, every element to 
the right of what you deleted has to shift over to the left by one unit. This takes O(n) time and 
on large arrays can be very costly. Because the next element has shifted over to the left, its index 
is now the current index, so you don't want to increment the index variable. 
In the implementation, the relocation logic is implemented as a ``swap and pop'' in which we swap the 
current element with the last one, then we delete the last element, which was the current element. 
This leads to no shifting of elements, and is logically equivalent (for our use case), even if the 
order of elements, relative to the deleted element, is no longer the same.
}

If we couldn't relocate the watched literal's index, we want to increment the index variable 
so that on the next iteration it will point to a new element. In the case that we couldn't 
relocate the watched literal, it means that every element in the clause, except for possibly the 
element pointed to by the second watched literal, evaluates to False. If the element pointed to by 
the second watched literal also evaluates to False, then we've encountered a contradiction in our 
assignment. However, if the second watched literal evaluates to True, then the clause is already 
satisfied and we need to do nothing except for move onto the next variable to propagate on. 
And if the second watched literal evaluates to Unassigned, then we can assign it the value it 
needs to evaluate to True and we then propagate on that newly assigned unit literal. 

If we encounter a contradiction, we must undo all the propagation we've done, including 
the initial decision before returning a null value to signify that the propagation failed. 
However, if it has succeeded, then we return the array of the variable indices that we've 
assigned.

So one of the big aspects of watched literals is how we get away with doing so little 
amount of work. And the answer is something called the ``watched literal invariant.''
This invariant is that at all times, at least one of the watched literals doesn't 
evaluate to False. I don't expect this to make everything click immediately, so let's 
explain more. A clause is only a contradiction when everything evaluates to False.
However, we don't want to check every single clause when we assign only one variable.
And even checking the negative of the variable we assigned is too much. So instead, we 
change which variables are watching which clauses every now and then. We pick at most 
two variables to watch a singular clause with, and each one watches their respective clauses
until they are assigned a value that makes them evaluate to False. In that case, there is 
the possibility that the watched literal invariant may be violated (ie, there is a possibility
that both watched literals may evaluate to False), but just to be sure, we try to change the 
variable watching that same clause. 

We scan the clause for any unwatched literal that is Unassigned or True, and if we find one, 
we've restored the watched literal invariant, so we don't need to propagate or do anything 
else for that matter. However, if we can't find one, then we're in serious danger of the 
watched literal invariant being violated. So our only hope is for the second watched literal 
to either already be True (we need to do nothing), or Unassigned (we assign it and propagate).

Thanks to that invariant, we also don't need to do any undoing when we backtrack other 
than undoing what we propagated. So long as the watched literal invariant holds true 
when we undo our assignments, the watched literals don't need to move, which means 
besides the undoing of what we propagated, we get zero cost backtracking.

\subsection{Why Two Watches?}
This now begs the question, why two watched literals? What is special about the number $2$?
Why can't we have just one watched literal? In order to answer this question, we need to ask 
ourselves: ``what can we do with two watched literals that we can't do with one?'' 
When we have two watched literals, we can efficiently check for whether or not we can 
propagate. How so?

After we try propagating on a literal, we scan the whole clause except for the indices 
of the first and second watched literals. Next, if everything in the clause evaluates to 
False, our only hope to propagate is on the second watched literal. So if that evaluates to 
False as well, we have found a contradiction, and if it doesn't we have found a literal we 
can propagate on.

But why is that more efficient that just one watched literal? Let's think of the worst case 
scenario for if we just had one watched literal. What does that look like? Suppose we were 
watching the first literal in the clause and every other literal in the clause evaluated 
to False. Ie, we need to propagate on this watched literal. How would it know it needs to 
propagate when it was never alerted by anything becoming False? It would have no idea 
it needs to do anything until that watched literal was falsified, and by that point,
everything in the clause would evaluate to False, so the clause would be unsatisfied,
only then would we backtrack, but when we backtrack, we'd still have no way of knowing if 
the clause is unit. 

So with only one watched literal, in the worst case scenario, which is that the literal
we're watching is the last literal in the clause to be falsified, we'd be doing the same 
amount of work as normal unit propagation, perhaps even more. 

Visually speaking, this is what I'm talking about:
\begin{center}
    \[ \blackinwhitesquare \square \square \]
    \[ U U U \]
\end{center}
The last literal becomes False, but we're not watching it so nothing moves.
\begin{center}
    \[ \blackinwhitesquare \square \square \]
    \[ U U F \]
\end{center}
Next the second literal becomes False, but we're not watching it so nothing happens.
\begin{center}
    \[ \blackinwhitesquare \square \square \]
    \[ U F F \]
\end{center}
Although we should perform unit propagation, our watched literal is completely blind 
to the fact that it is the last Unassigned literal in the clause, so it remains as is.
If luck has it, it could become True and all would be well, but this is reality, so more 
likely than not, it will also become False.
\begin{center}
    \[ \blackinwhitesquare \square \square \]
    \[ F F F \]
\end{center}
Now we scan the clause, and we see that we can't relocate the watched literal, so we 
backtrack on our assignment.
\begin{center}
    \[ \blackinwhitesquare \square \square \]
    \[ U F F \]
\end{center}
Since the watched literal was never aware if both of the other literals became False 
in one pass of the algorithm or if it was across multiple decisions, it still doesn't 
know that it's the last Unassigned literal. It has to scan the clause each time to know 
whether or not it's the last Unassigned literal, which is no better than regular unit
propagation.

\subsection{Watched Literals - Results}

After implementing watched literals, our runtime decreases dramatically. 
For uf20.91, solving all 1000 instances takes \textbf{349ms}, as opposed to the 
\textbf{1548ms} from pure literal elimination and the \textbf{868ms} from normal unit propagation.
And for UUF50.218.1000, it takes only \textbf{7130ms} as opposed to the \textbf{231,107ms} 
from pure literal elimination and the \textbf{140,411ms} from normal unit propagation.
Watched literals took us down from minutes to seconds for the unsatisfiable 
problems, which is amazing for us.

\begin{FVerbatim}
On all 1000 equations in uf20-91:

Watched Literals:          349ms (~0.35 seconds).
Pure Literal Elimination:  1548ms (~1.55 seconds).
Unit Propagation:          868ms (~0.87 seconds).
Recursive Backtracking:    1,540,613ms (~25.68 minutes).
Brute Force:               1,770,851ms (~29.51 minutes).
Short-Circuited (Basic):   49,874ms (~50 seconds).
Parallel:                  10,589ms (~10.5 seconds).

On all 1000 equations in UUF50.218.1000:

Watched Literals:         7,130ms (~7.13 seconds).
Pure Literal Elimination: 231,107ms (~3.85 minutes).
Unit Propagation:         140,041ms (~2.34 minutes).
\end{FVerbatim}

\section{Iterative DPLL}

Now that we've implemented DPLL to the best of our ability, it's time to 
take it one step further. We will now implement it iteratively rather than 
recursively. Why? Remember how I said in the beginning that DPLL is the foundation 
for modern SAT solvers?
\footnote{I can't remember if I actually said that}
Well iterative DPLL is the foundation for the next chapter's algorithm: CDCL.
Implementing iterative DPLL is imperative in order to create a modern SAT solver
worth its salt.
\footnote{Honestly, ``Salt'' is a better name than ``Saturn.'' I might change it.}

As said in the beginning of the chapter, DPLL forces us to traverse the decision 
tree as well, a tree. And that is still true, even in iterative DPLL. What changes 
now is that instead of an implicit path, we now have an explicit one. Suppose we 
are on the first decision level, so we decide variable $1 = T$, and that forces 
$3 = F$, $5 = F$, and $6 = T$. In recursive DPLL, we assign these variables into 
the \textit{assignment} or \textit{vars} field, but we don't really use them until 
we get to the decision level that matches their index or the evaluation phase. 

With iterative DPLL, there are still decision levels, but the representation is 
now very, very different bcause instead of a decision \textit{tree}, we now use a 
decision \textit{trail}. The trail is reminiscent of a queue in the sense that 
we only operate at the end of the data structure. It contains some information 
that we need in order to continue mimicking the behavior of recursive DPLL. Each 
trail entry contains the literal that was forced, the decision level it's a part of,
and the clause that forced that decision. We append these trail entries every single 
time we force a variable's assignment. The next question is how do we know which 
variables were decided for and which variables were inferred from propagation?
Well, for decisions, there is no reason clause because there was no reason to 
begin with other than ``because I said so.'' Therefore, the variables that were 
decided for don't get a reason clause, but the ones that were inferred do. This 
is great because it makes it easy to tell where the beginning of one decision 
starts and where it ends.

However, it would be nice to have some way to keep track of where the decisions
are in the trail other than scanning for entries without a reason clause. Because 
of that, we will also keep another trail specifically for the decision markers.
This is also helpful because then we can store whether or not we've ``tried the 
other branch'' when we backtrack. 

While on the topic of backtracking, it would be nice to discuss it some more as well. 
Since we can't simply go up the call stack and undo variable assignments, we have to 
do it in the iterative way. When we backtrack, we first backtrack up to the decision
entry. This is done by popping off the trail and undoing the propagated assignments. 
Once that's done, you then check if you've ``tried the other branch'', and if you 
haven't, you flip the decision literal's assignment and then retry. If you have 
already tried the other branch, you pop off the decision entry, and then repeat 
this process to the next decision entry. This process continues until you either 
find a decision that has not been switched yet, or you run out of decision 
entries.

With most of the exposition out of the way, I think it would be good to 
look at the changes we will make to our interface and algorithms, starting 
with our SAT solver.

\begin{lstlisting}
    public interface:
        SATSolver(string input)
        bool solveCNF()

    private interface:
        enum var
        {
            False
            True
            Unassigned
        }

        TrailEntry {
            int lit,
            int clauseIdx
        }

        DecisionLevel {
            int idx,
            bool triedOtherBranch
        }

        void parse(string equation)
        bool propagate()
        bool createWatched()

        int numVars
        int numClauses
        int trailHead

        int[][] clauses
        bool[] satisfied
        
        TrailEntry[] trail
        DecisionLevel[] trailStarts

        optional var[] vars
        map<int, int[]> var_to_clause
        map<int, (int, int)> clause_to_vars
\end{lstlisting}

The first notable change is that \textit{propagate()} now returns a boolean instead 
of a delta. This is because it will no longer be responsible for undoing propagation 
in response to a contradiction, thus achieving zero-cost backtracking. Next, 
we now have a \textit{trail} and a \textit{trailHead}, which points to the last 
made decision entry. \textit{trailStarts} keeps track of all the indices of the 
decision entries. One last thing to note is that the \textit{TrailEntry} data 
structure uses indices for clauses instead of the actual clause itself. 
There's a very good reason for this, and I will reveal the answer in the next chapter,
but until then, why do you think it might be that way? Hint: it's not because 
indices are easy for accessing. 

The next change we need to make is to our algorithm to create watched literals. 
It's the same for basically $90\%$ of the algorithm, except for the unit clause 
case, which I'll highlight.

\begin{lstlisting}
    bool createWatched(int[] assignment) {
        for (i = 0 up to numClauses) {
            clause = clauses[i]
            ...
            else if (clause.size == 1) {
                int var = clause[0]
                ...
                append(trail, {var, i})
                ...
            }
            ...
        }
        return True
    }
\end{lstlisting}

Rather than enqueuing the variable, we create a new trail entry for each unit 
clause. These ones are all part of the \textbf{root} decision level, which we 
mentally reserve for the 0th decision level. If we backtrack up to this level (if it exists)
or before it, we've encountered a root level contradiction, so the equation is 
unsatisfiable. Because not every equation has unit clauses, it's important to 
remember that this decision level may not always exist, and so we should be 
keep that in mind for our solving algorithm. 

The next algorithm we're going to change is the \textit{propagate()} algorithm 
because the changes are also very minimal.

\begin{lstlisting}
    bool propagate(int[] assignment) {
        while (trailHead < trail.size) {
            int prop = trail[trailHead].lit
            trailHead = trailHead + 1
            ...
            for (i = 0 up to watched_clauses.size) {
                clause_idx = watched_clauses[i]
                int[] clause = clauses[clause_idx]
                (int, int) watches = clause_to_var[clause_idx]
                ...
                if (relocated == False) {
                    ...
                    int other_watch = clause[watches.second];
                    int idx = abs(other_watch)
                    if (assignment[idx] == Unassigned) {
                        ...
                        append(trail, {other_watch, clause_idx})
                    } else {
                        return false
                    }
                }
            }
        }
        return true
    }
\end{lstlisting}

The main difference is that we just \textit{trailHead} to iterate through the trail, and 
instead of enqueuing, we just append to the trail. We also return True/False depending on 
whether or not propagation succeeded.

Lastly, we look at our \textit{solve()} method, which has the most notable changes.
There are two ways to implement this, depending on what makes sense to you. I call 
them ``assignment before propagation'' and ``propagation before assignment.'' I
call them these because that's what they are. Before showing the pseudocode, I 
want to talk about the algorithm a bit more. The beginning is mostly the same:
create our watched literals, check that it succeeded, return False if it didn't. 
Otherwise, we do a root level propagation and then check if it succeeds as well.
However, the main confusion comes from handling the case of when there is a 
0th decision level and when there isn't. There are multiple ways to go about this.
After the beginning, we enter the solving loop. 

There are two main parts: assigning and backtracking, with a call to the propagate 
method separating them. Assigning is now scanning for the first Unassigned variable,
forcing it True, creating a new trail entry, and then updating the trail head. I'm 
sure there's a better way to keep track of the next unassigned variable other than 
scanning the entire assignment list on each loop iteration, however, for simplicity's 
sake, we're implementing it the naive way. 

Since we've already discussed backtracking, we will now get into the algorithm 
itself:

\begin{lstlisting}
    bool solve() {
        var[1..numVars] assignment

        for each (var in assignment) {
            var = Unassigned
        }

        if (createWatched(assignment) == False) {
            return False
        }

        if (trail.size > 0) {
            append(trailStarts, {0, true})

            if (propagate(assignment, 0) == False) {
                return False
            }
        }

        while (True) {
            bool succeeded = propagate(assignment)

            if (!succeeded) {
                while (trailStarts.size > 0) {
                    DecisionLevel decision = trailStarts.last

                    while (trail.size > decision.idx + 1) {
                        int lit = trail.last.lit
                        assignment[abs(lit)] = Unassigned
                        erase(trail, trail.last)
                    }

                    if (decision.triedFalse == False) {
                        TrailEntry entry = trail.last

                        entry.lit = -entry.lit
                        assignment[abs(entry.lit)] = False
                        decision.triedFalse = True

                        trailHead = decision.idx
                        break
                    } else {
                        int lit = trail.last.lit
                        assignment[abs(lit)] = Unassigned
                        erase(trail.last)
                        erase(trailStarts.last)
                    }
                }

                if (trailStarts.size == 0) {
                    return False
                }

                continue
            }
        }

        //Assignment code
        int firstUnassigned = -1
        for (i = 1 up to numVars) {
            if (assignment[i] == Unassigned) {
                firstUnassigned = i
                break
            }
        }

        if (firstUnassigned == -1) {
            return True
        }

        assignment[firstUnassigned] = True
        append(trail, {firstUnassigned, -1})
        append(trailStarts, {trail.size - 1, False})
        trailHead = trailStarts.last.idx
    }
\end{lstlisting}

Here is the pseudocode for solving, written as ``propagate then assign." Quick question:
when there's no 0th decision level and we propagate, how do we know that none of the 
propagated variables will have the wrong decision level? Ie, how do we know that they'll 
have a decision level of $0$ instead of $1$, since we haven't made any decisions yet?
This is kind of a trick question, the answer is that if there is no 0th level,
\textit{propagate()} will do nothing because the while-loop condition will immediately 
return False, and so the method will terminate, assigning nothing! 

Instead of showing you the code for ``assign then propagate'' (it is NOT just taking 
the assigning code and moving it up before the call to \textit{propagate()}), let's 
run through how this works. When there's a 0th level, we obviously want to mark it 
in our \textit{trailStart}, but why do we want to mark that the ``other branch'' has 
already been tried? Well that's because those assigned by the root level propagation 
are forced into those values, so it doesn't make sense to swap their values to be 
their opposites since it will immediately be an unsatisfiable assignment anyways. 

Next, in the main solving loop, the backtracking code is just ``grab the last 
decision index, undo until the root decision, if it's already tried False, undo 
the decision entirely and repeat with the previous one, and if not, try False 
and continue.''

Lastly, let's talk about doing assign then propagation. Here's why it doesn't
make total sense to start out with propagation:

\begin{itemize}
    \item we do a root level propagation if there needs to be one.
    \item if there is no root level propagation to be done, it still
            doesn't make sense to immediately propagate since there's 
            nothing to propagate.
    \item if there was a root level propagation, it doesn't make sense 
            to do another one immediately since we've just completed one.
\end{itemize}

So if you wanted to do assign and then propagate, how would that be done? 
Earlier, I said that it's not just taking the code for assigning and then 
moving it before propagation, but why? Well, suppose we fail a propagation. 
We backtrack, and then suppose we haven't tried the other branch for that 
decision. We encounter a \textit{continue} statement and we go back to the 
top. If we simply move the assigning code up to the top of the loop,
what's the issue? Well remember, we just flipped a decision, so we should
actually be \textbf{\textit{propagating}} instead of assigning. If we 
assign immediately after flipping the decision, we now have two decisions 
at the same time, and that's bad because what if they are contradictory to 
one another, ie one would force a variable False and the other True? 
Unfortunately, the propagation code wouldn't catch that because it assumes 
everything in the trail is valid, but we just made the trail invalid. So 
to the propagation code, instead of it catching that a variable is forced to be 
False, but it's True, or vice versa, it may just ignore it because it focuses 
mainly on Unassigned and False literals.

So if we want to do assign then propagate, how do we make sure that 
we don't accidentally push something onto the trail when we shouldn't have?
We need to wrap the assigning code in a conditional statement that checks that 
\textit{trailHead == trail.size}, if it doesn't, then that means that there's 
something to be propagated on because we're not at the end of the trail, but 
if it does, then we've reached the end, so we should check if we've assigned 
everything and can be done, or if we need to assign something new.

\subsection{Iterative DPLL - Results}
Here are the results after this huge refactor:

(Assignment before propagation)\\
uf20.91: \textbf{328ms}, down 21ms from \textbf{349ms} by Watched Literals\\
UUF50.218.1000: \textbf{6369ms}, down 761ms from \textbf{7130ms} by Watched Literals

(Propagation before assignment)\\
uf20.91: \textbf{324ms}\\
UUF50.218.100: \textbf{6365ms}

\begin{FVerbatim}
On all 1000 equations in uf20-91:

Iterative DPLL (AbP):      328ms (~0.33 seconds).
Iterative DPLL (PbA):      324ms (~0.32 seconds).
Watched Literals:          349ms (~0.35 seconds).
Pure Literal Elimination:  1548ms (~1.55 seconds).
Unit Propagation:          868ms (~0.87 seconds).
Recursive Backtracking:    1,540,613ms (~25.68 minutes).
Brute Force:               1,770,851ms (~29.51 minutes).
Short-Circuited (Basic):   49,874ms (~50 seconds).
Parallel:                  10,589ms (~10.5 seconds).

On all 1000 equations in UUF50.218.1000:

Iterative DPLL (AbP):     6,369ms (~6.37 seconds).
Iterative DPLL (PbA):     6,365ms (~6.37 seconds).
Watched Literals:         7,130ms (~7.13 seconds).
Pure Literal Elimination: 231,107ms (~3.85 minutes).
Unit Propagation:         140,041ms (~2.34 minutes).
\end{FVerbatim}

Should you do propagation before assignment simply because it's faster? It 
really doesn't matter, you do you. However, for the rest of this book, we will be 
using assignment before propagation because it's easier for me to understand, and it 
follows online pseudocode examples better than propagate before assignment.

The speed gain on satisfiable problems is very miniscule, however, if you ever want 
proof that iterative algorithms are faster than recursive ones, the unsatisfiable 
problem time decrease is pretty good evidence.

\end{document}